% FC'27 SoK camera-ready: 20 pages + references and appendices, LNCS, named author block.
% Ported from the AsiaCCS 2027 submission (src_driver_acm.latex) on 2026-09-20.
% Compile: pdflatex main ; bibtex main ; pdflatex main ; pdflatex main
\documentclass[runningheads]{llncs}

\usepackage[T1]{fontenc}
\usepackage{amsmath,amssymb}
\usepackage{booktabs}
\usepackage{graphicx}
\usepackage{xcolor}
\usepackage{enumitem}
\usepackage{multirow}
\usepackage{array}
\usepackage{makecell}
\usepackage{threeparttable}
\usepackage{tikz}
\usetikzlibrary{positioning,shapes.geometric,arrows.meta,calc,fit,backgrounds,shapes.misc}
\usepackage[hidelinks]{hyperref}
\usepackage{xurl}

% Typesetting robustness: let TeX stretch to avoid overfull lines with long identifiers
\setlength{\emergencystretch}{2em}
\renewcommand{\topfraction}{0.9}
\renewcommand{\floatpagefraction}{0.8}

% acmart compatibility: \Description carries figure alt text; llncs has no
% equivalent, so absorb it (optional short form + mandatory long form).
\newcommand{\Description}[2][]{}

% ---- Convenience macros (ported unchanged) ----
\newcommand{\coldcard}{\textsc{ColdCard}}
\newcommand{\dualec}{\textsc{Dual\_EC\_DRBG}}
\newcommand{\stage}[1]{\textsf{S#1}}
\newcommand{\bits}[1]{$2^{#1}$}
\newcommand{\etal}{et al.\@}
\newcommand{\ie}{i.e.,\@}
\newcommand{\eg}{e.g.,\@}
\newcommand{\prevent}{\textbf{P}}
\newcommand{\detect}{\textbf{D}}
\newcommand{\contain}{\textbf{C}}
% Harvey balls for the remedy-coverage matrix (Table~\ref{tab:coverage}).
\DeclareRobustCommand{\hbfull}{\tikz[baseline=-0.55ex]\fill (0,0) circle (0.42ex);}
\DeclareRobustCommand{\hbhalf}{\tikz[baseline=-0.55ex]{\draw[line width=0.13ex] (0,0) circle (0.42ex);
  \fill (0,0.42ex) arc (90:270:0.42ex) -- cycle;}}
\DeclareRobustCommand{\hbnone}{\tikz[baseline=-0.55ex]\draw[line width=0.13ex] (0,0) circle (0.42ex);}
\DeclareRobustCommand{\hbna}{\textcolor{black!35}{--}}

\begin{document}

% FC'27 requires SoK titles to begin with the literal text "SoK:".
\title{SoK: Randomness Failures and Their Remedies in Deployed Cryptographic Systems}

\titlerunning{SoK: Randomness Failures and Their Remedies}
\author{Devharsh Trivedi\orcidID{0000-0001-6374-7249}}
\authorrunning{D. Trivedi}
\institute{Mathematics, Engineering and Computer Science,\\
LaGuardia Community College, City University of New York,\\
Long Island City, NY, USA\\
\email{devharsh.trivedi86@lagcc.cuny.edu}}

\maketitle

\begin{abstract}
In July 2026 attackers drained wallets whose seeds had been generated on \coldcard{} hardware
devices, with published losses between about \$88 million and \$130 million. A build error had linked
seed generation to a non-cryptographic software generator for more than five years. Published
analyses disagree on how much entropy the seeds kept, from about 40 or 72 bits by the vendor's
estimate to \bits{32} distinguishable outputs under Block's stronger assumptions, but all fall far
short of 128 bits. The incident is the latest in a thirty-year record of randomness failures in
TLS and SSH hosts, smart cards, virtual machines, processors, and wallets. This Systematization
of Knowledge adopts the generator architecture of NIST SP 800-90 for the first five stages of a
randomness pipeline and adds two loci the standards leave out: the consumption of random values
by protocols, and the build and integration layer that decides which generator is actually
linked. We code twenty-eight deployed incidents on this model and cross them with fourteen remedy
families in a coverage matrix built by stated and released rules. Read by column, the matrix
shows that statistical testing and certification prevent none of the failures, which confirms
in the field a limit long known in principle, and that only build-path verification reaches the
build and integration layer. Read by row, it shows that for 22 of the 24 preventable failures a
preventing remedy was at least partly available when the system shipped. We close with open
research questions tied to the empty parts of the matrix and a prioritized checklist for
implementers.

\keywords{Random number generation \and Entropy \and CSPRNG \and DRBG \and ECDSA nonces \and Hardware wallets \and Cryptographic engineering.}
\end{abstract}

\input{sections/01-intro.latex}
\input{sections/02-background.latex}
\input{sections/03-taxonomy.latex}
\input{sections/04-incidents.latex}
\input{sections/05-remedies.latex}
\input{sections/06-effectiveness.latex}
\input{sections/06b-related.latex}
\input{sections/07-discussion.latex}
\input{sections/08-conclusion.latex}

\begin{credits}
\subsubsection{\ackname} The author thanks the anonymous FC 2027 reviewers, whose comments
shaped this version. The author used an AI assistant (Claude, Anthropic) for editing, reference
checking, figure drafting, and the artifact scripts, and reviewed all content and takes
responsibility for it.

\subsubsection{\discintname} The author has no competing interests to declare that are relevant
to the content of this article.
\end{credits}

\bibliographystyle{splncs04}
\bibliography{references}

\appendix
\input{sections/09-appendix.latex}

\end{document}
